\section{The Stripped-Down Model}
\label{sec:stripped}

The preceding sections establish, separately, that the lag basis is a
non-lever, that estimated geometry loses to fixed geometry everywhere it was
tried, and that the exogenous operator's fitted structure is concentrated in a
handful of directions. Taken together those results have a constructive
consequence that none of them states on its own: almost all of the apparatus
can be deleted.

\subsection{Fifty-three columns}
\label{sec:dumb_design}

The design is:

\begin{itemize}
  \item the twelve base-2 HAR rungs on $RV$ (Section~\ref{sec:descriptive});
  \item one number per exogenous channel---its twelve-rung ladder contracted
    against a fixed decaying weight,
    \[
      a_c(t) \;=\; \sum_{k} u_k^{-1}\, Z_{c,k}(t),
      \qquad u_k \in \{1, 2, 4, \dots, 2048\},
    \]
    which replaces $41 \times 12 = 492$ columns with $41$;
  \item ridge, one penalty, chosen on a validation region that precedes the
    evaluation region under the same rolling protocol;
  \item a rolling $24{,}000$-bar training window refit every $1{,}000$ bars.
\end{itemize}

Nothing else. No embedding, no nearest neighbours, no principal components, no
singular value decomposition, no tensor factorisation, no manifold, no
model-class search.

The weight $u_k^{-1}$ is not fitted. It is the shape the operator's own
directional analysis converged on (Section~\ref{sec:knn_metric_race}), and
Section~\ref{sec:dumb_results} shows the result is insensitive to the exponent
within the range tested but not outside it.

\subsection{It beats the elaborate version}
\label{sec:dumb_results}

Table~\ref{tab:dumb} runs four designs over the full walk-forward: identical
panel, identical rows, identical rolling protocol, penalties chosen the same
way.

\begin{table}[H]
\centering
\caption{Full walk-forward, $218{,}909$ bars, rolling $24{,}000$-bar window at
cadence $1{,}000$. The last column is matched-row Diebold--Mariano against the
stripped-down design; negative favours it.}
\label{tab:dumb}
\begin{tabular}{lrrrl}
\toprule
Design & Columns & QLIKE & MSE & DM vs.\ stripped \\
\midrule
backbone only            &  12 & 0.13571 & $4.166 \times 10^{-11}$ & $-0.00291$ ($t = -9.59$) \\
\textbf{stripped-down}   &  53 & \textbf{0.13280} & $4.389 \times 10^{-11}$ & --- \\
ridge on 492 raw columns & 504 & 0.13617 & $4.364 \times 10^{-11}$ & $-0.00338$ ($t = -2.18$) \\
shape-projection penalty & 504 & 0.13415 & $4.541 \times 10^{-11}$ & $-0.00135$ ($t = -5.63$) \\
\bottomrule
\end{tabular}
\end{table}

Three readings.

\paragraph{The elaborate estimator is a feature in disguise.} The
shape-projection penalty on $504$ columns and the stripped-down design on $53$
are not close substitutes---they are the same estimator computed two ways. The
effective-degrees-of-freedom accounting says so directly: the penalty spends
$\approx 54$ degrees of freedom on its $505$-column design, which is
$1 + 12 + 41$, the dimension of its own null space
(Section~\ref{sec:effective_df}). Precomputing the $41$ amplitudes and applying
plain ridge does not approximate that, it \emph{implements} it, and does so
slightly better ($-0.00135$, $t = -5.63$) because it does not also have to
estimate $492$ coefficients it will then shrink away. On a separate evaluation
window the two agreed to one part in $10^{5}$ ($0.11783$ against $0.11784$).

\paragraph{The raw parameterisation is actively harmful.} Plain ridge on the
$492$ raw ladder columns is the \emph{worst} design in the table, $0.00046$
behind the twelve-column backbone it is supposed to improve. Carrying the block
in its natural form does not merely waste degrees of freedom; it loses to not
carrying it at all.

\paragraph{The gain is small and the robustness is not.} The stripped-down
design improves on the backbone by $0.00291$ ($t = -9.59$), which is the same
order as every other information effect in this paper. What is not the same
order is its behaviour in the one era where the exogenous block fails. In
2022-03 to 2024-04, ridge on raw columns loses to the backbone by $+0.03321$;
the stripped-down design loses by $+0.00347$, a tenfold containment, and the
bespoke penalty by $+0.00666$. An amplitude averages a channel's whole rung
ladder, so a single detonating rung is diluted rather than carried at full
weight---which is why the summary is more robust than the thing it summarises.

\subsection{Why the two are the same estimator}
\label{sec:effective_df}

A quadratic penalty and a kernel are the same object: minimising
$\lVert y - X\beta \rVert^2 + \beta' P \beta$ is kernel ridge with
$k(x, x') = x' P^{-1} x'$, an identity we verify numerically to $3.1 \times
10^{-15}$ relative. The kernel view supplies the honest measure of how much
freedom a fit uses,
\[
  \mathrm{df}(\lambda) \;=\; \operatorname{tr}\!\big( (G + \lambda P)^{-1} G \big),
\]
which is not the column count but how the penalty meets the data spectrum.

\begin{table}[H]
\centering
\caption{Effective degrees of freedom, rolling $24{,}000$-bar windows. Both
penalised designs carry the same $505$ columns.}
\label{tab:df}
\begin{tabular}{lrrr}
\toprule
Era & backbone & ridge & shape-projection \\
\midrule
2007--2009 & 13.0 & 93.0 & 47.2 \\
2013--2015 & 13.0 & 98.9 & 55.6 \\
2020--2022 & 13.0 & 98.6 & 56.3 \\
2022--2024 & 13.0 & 104.7 & 60.0 \\
\bottomrule
\end{tabular}
\end{table}

The shape-projection penalty spends $\approx 54$ degrees of freedom---
$1 + 12 + 41$, precisely the dimension of its null space, one amplitude per
channel. That is the arithmetic behind Section~\ref{sec:dumb_results}: the
penalty was never using its $492$ columns, so precomputing what it does use
loses nothing.

Splitting that freedom by Gram eigenvalue block makes the dense-but-weak
finding quantitative rather than rhetorical. The design's first principal
component carries $96.2\%$ of the variance and $0.02\%$ of the explained
target, with the predictable variation in components $6$--$20$
($0.18\%$ of variance, $79.9\%$ of explained target). Ridge spends $72.08$ of
its $\approx 95$ degrees of freedom in components $21$ and beyond and only
$14.12$ in the band that predicts; the shape-projection penalty spends
$0.30$, $0.13$ and $0.07$ across the three blocks. Ridge preserves exactly the
directions carrying no signal, because it shrinks direction $i$ by
$d_i / (d_i + \lambda)$ and those are the directions with the largest $d_i$.

\subsection{What this costs to adopt}

The stripped-down design has one hyperparameter, no estimated geometry, and a
feature construction expressible in one line. Against that, this paper's
earlier apparatus required a $41 \times 12$ operator, its singular value
decomposition, session-conditional frames, CP and Tucker factorisations,
rank-constrained priors and a bespoke projection penalty, and reached a worse
number.

We report this as the section's finding rather than as an appendix. The
methodological claim of Section~\ref{sec:tuning_nonlever}---that in the
dense-but-weak regime complexity is a cost to be justified---applies to our own
construction, and applied to it, most of it does not survive.
